
\textup{B)} We now use the associativity of \(G\) to write
\[
  p_\alpha(u)p_{-\alpha}(v)p_{-\alpha}(w)
    =p_\alpha(u)p_{-\alpha}(v+w).
\]
There is an open subset \(L\) of
\(\mathbb G_{a,S}^{\,3}\), containing the zero section, such that
\((u,v,w)\in L(S')\) is equivalent to
\[
\begin{gathered}
  (u,v)\in V(S'),\qquad
  \bigl(uE(uv)^{-1},w\bigr)\in V(S'),\qquad
  (u,v+w)\in V(S'),\\
  uv\in W(S'),\qquad
  uwE(uv)^{-1}\in W(S'),\qquad
  u(v+w)\in W(S').
\end{gathered}
\]
Using formula \((+)\), we immediately obtain, for
\((u,v,w)\in L(S')\),
\SGARefTarget{sga3:xx:2.2:1}
\begin{equation}\tag{1}
  E(uv+uw)=E\bigl(uwE(uv)^{-1}\bigr)E(uv),
\end{equation}
\SGARefTarget{sga3:xx:2.2:2}
\begin{equation}\tag{2}
  H(uv+uw)=H\bigl(uwE(uv)^{-1}\bigr)H(uv),
\end{equation}
and
\SGARefTarget{sga3:xx:2.2:3}
\begin{equation}\tag{3}
\begin{split}
  (v+w)F(uv+uw)^{-1}
   ={}&\alpha\bigl(H(uv)^{-1}\bigr)\,
        wF\bigl(uwE(uv)^{-1}\bigr)^{-1}\\
     &{}+vF(uv)^{-1}.
\end{split}
\end{equation}

It follows immediately from the definition of \(L\) that
\((1,0,0)\in L(S)\). Consider therefore
\[
  L\cap_T
    \bigl(1\times_S\mathbb G_{a,S}\times_S\mathbb G_{a,S}\bigr)
      =1\times_S M.
\]
The scheme \(M\) is an open subset of
\(\mathbb G_{a,S}^{\,2}\) containing the section \((0,0)\).
For \((v,w)\in M(S')\), the elements
\[
  v,\qquad wE(v)^{-1},\qquad v+w
\]
belong to \(W(S')\), and
\SGARefTarget{sga3:xx:2.2:1-prime}
\begin{equation}\tag{1'}
  E(v+w)=E\bigl(wE(v)^{-1}\bigr)E(v),
\end{equation}
\SGARefTarget{sga3:xx:2.2:2-prime}
\begin{equation}\tag{2'}
  H(v+w)=H\bigl(wE(v)^{-1}\bigr)H(v),
\end{equation}
\SGARefTarget{sga3:xx:2.2:3-prime}
\begin{equation}\tag{3'}
\begin{split}
  (v+w)F(v+w)^{-1}
   ={}&\alpha\bigl(H(v)\bigr)^{-1}
        wF\bigl(wE(v)^{-1}\bigr)^{-1}\\
     &{}+vF(v)^{-1}.
\end{split}
\end{equation}

Finally, consider the morphism from \(M\) to
\(\mathbb G_{a,S}^{\,2}\) defined set-theoretically by
\[
  (v,w)\longmapsto\bigl(v,wE(v)^{-1}\bigr).
\]
\footnote{The editors corrected what follows.}
It preserves the section \((0,0)\) and induces an isomorphism of
\(M\) onto an open subset \(N\) of
\(\mathbb G_{a,S}^{\,2}\) containing the zero section; the inverse
isomorphism is
\[
  (x,y)\longmapsto\bigl(x,yE(x)\bigr).
\]
\footnote{Thus the editors made the change of variables
\(x=v,\ y=wE(v)^{-1}\), or equivalently
\(v=x,\ w=yE(x)\).}
We have therefore proved the following assertion:

There is an open subset \(N\) of
\(\mathbb G_{a,S}^{\,2}\), containing the zero section, such that,
if \((x,y)\in N(S')\), then \(x\), \(y\), and
\(x+yE(x)\)\footnote{The editors corrected \(yE(x)\) to
\(x+yE(x)\).} belong to \(W(S')\), and
\SGARefTarget{sga3:xx:2.2:1-double-prime}
\begin{equation}\tag{1''}
  E\bigl(x+yE(x)\bigr)=E(x)E(y),
\end{equation}
\SGARefTarget{sga3:xx:2.2:2-double-prime}
\begin{equation}\tag{2''}
  H\bigl(x+yE(x)\bigr)=H(x)H(y),
\end{equation}
\SGARefTarget{sga3:xx:2.2:3-double-prime}
\begin{equation}\tag{3''}
\begin{split}
  \bigl(x+yE(x)\bigr)F\bigl(x+yE(x)\bigr)^{-1}
   ={}&xF(x)^{-1}\\
     &{}+\alpha\bigl(H(x)\bigr)^{-1}yE(x)F(y)^{-1}.
\end{split}
\end{equation}

\textup{C)} Arguing in the same way with left associativity, we
obtain the following assertion.\footnote{That is, one writes the
equalities resulting from
\(p_\alpha(t)p_\alpha(u)p_{-\alpha}(v)
  =p_\alpha(t+u)p_{-\alpha}(v)\), and sets
\(v=1,\ x=u,\ t=yF(u)\), or equivalently
\(y=tF(u)^{-1}\).}

There is an open subset \(N'\) of
\(\mathbb G_{a,S}^{\,2}\), containing the zero section, such that,
if \((x,y)\in N'(S')\), then \(x\), \(y\), and
\(x+yF(x)\)\footnote{The editors corrected \(yF(x)\) to
\(x+yF(x)\).} belong to \(W(S')\), and
\SGARefTarget{sga3:xx:2.2:4-double-prime}
\begin{equation}\tag{4''}
  F\bigl(x+yF(x)\bigr)=F(x)F(y),
\end{equation}
\SGARefTarget{sga3:xx:2.2:5-double-prime}
\begin{equation}\tag{5''}
  H\bigl(x+yF(x)\bigr)=H(x)H(y),
\end{equation}
\SGARefTarget{sga3:xx:2.2:6-double-prime}
\begin{equation}\tag{6''}
\begin{split}
  \bigl(x+yF(x)\bigr)E\bigl(x+yF(x)\bigr)^{-1}
   ={}&xE(x)^{-1}\\
     &{}+\alpha\bigl(H(x)\bigr)^{-1}yF(x)E(y)^{-1}.
\end{split}
\end{equation}

We are thus led to solve the ``functional equation'' \((1'')\).

\paragraph{\textbf{Lemma 2.3.}}
\SGARefTarget{sga3:xx:2.3}
\emph{Let \(S\) be a scheme, let \(W\) be an open subset of
\(\mathbb G_{a,S}\) containing the zero section, and let}
\[
  F\colon W\longrightarrow\mathbb G_{m,S}
\]
\emph{be an \(S\)-morphism. Suppose that \(F(0)=1\) and that there
is an open subset \(N\) of \(\mathbb G_{a,S}^{\,2}\), containing
the zero section, such that, for \((x,y)\in N(S')\), the elements
\(x\), \(y\), and \(x+yF(x)\) belong to \(W(S')\), and}
\SGARefTarget{sga3:xx:2.3:dagger}
\begin{equation}\tag{\(\dagger\)}
  F\bigl(x+yF(x)\bigr)=F(x)F(y).
\end{equation}
\begin{enumerate}
\item[\textup{\SGAInlineRefTarget{sga3:xx:2.3:item:i}(i)}] \emph{If \(S=\operatorname{Spec}(k)\) for a
  field \(k\), there is an \(a\in k\) such that
  \(F(x)=1+ax\).}
\item[\textup{\SGAInlineRefTarget{sga3:xx:2.3:item:ii}(ii)}] \emph{If
  \(a=F'(0)\in\Gamma(S,\mathscr O_S)\) is invertible, then
  \(F(x)=1+ax\).}
\end{enumerate}

\medskip
\noindent\emph{Proof.}
Under the hypotheses, we may differentiate the given equation at
\(x=0\), respectively at \(y=0\), and obtain
\SGARefTarget{sga3:xx:2.3:star}
\begin{equation}\tag{\(\ast\)}
  F'(y)\bigl(1+yF'(0)\bigr)=F'(0)F(y)
  \qquad\text{for }(0,y)\in N(S'),
\end{equation}
respectively
\[
  F'(x)F(x)=F(x)F'(0)
  \qquad\text{for }(x,0)\in N(S').
\]
Since \(F\) takes its values in \(\mathbb G_m\), the second
relation gives
\SGARefTarget{sga3:xx:2.3:star-prime}
\begin{equation}\tag{\(\ast'\)}
  F'(x)=F'(0)\qquad\text{for }(x,0)\in N(S').
\end{equation}
The first relation therefore gives
\[
  F'(0)\bigl(1+yF'(0)\bigr)=F'(0)F(y)
  \qquad\text{for }(y,0),(0,y)\in N(S').
\]
If \(a=F'(0)\) is invertible, it follows that
\[
  F(y)=1+ay
\]
for \(y\) a section of an open subset of \(W\) containing the zero
section, hence schematically dense in \(W\). This proves
\textup{\SGARefLink{sga3:xx:2.3:item:ii}{(ii)}}, and also proves \textup{\SGARefLink{sga3:xx:2.3:item:i}{(i)}} when \(F'(0)\ne0\).

Suppose \(F'(0)=0\). Then \((\ast')\) shows that \(F'(x)=0\)
for \(x\) near \(0\), and hence \(F'=0\) by schematic density.
If \(k\) has characteristic zero, \(F\) is a rational function
with zero derivative, and is therefore constant and equal to
\(F(0)=1\).

Suppose now that \(k\) has characteristic \(p\), and that \(F\)
is not constant.\footnote{The editors corrected what follows.}
There are an integer \(n>0\) and a rational function
\(F_1\in k(X)\) such that \(F_1'(X)\ne0\) and
\[
  F(X)=F_1\bigl(X^{p^n}\bigr)=F_1(X)^{p^n}.
\]
Substitution in the functional equation gives
\SGARefTarget{sga3:xx:2.3:dagger-one}
\begin{equation}\tag{\(\dagger_1\)}
  F_1\bigl(x+yF_1(x)^{p^n}\bigr)=F_1(x)F_1(y).
\end{equation}
Differentiating at \(x=0\) gives
\SGARefTarget{sga3:xx:2.3:star-one}
\begin{equation}\tag{\(\ast_1\)}
  F_1'(y)=F_1'(0)F_1(y),
\end{equation}
whereas differentiation of \((\dagger_1)\) at \(y=0\) gives
\SGARefTarget{sga3:xx:2.3:star-prime-one}
\begin{equation}\tag{\(\ast'_1\)}
  F_1'(x)F_1(x)^{p^n}=F_1(x)F_1'(0).
\end{equation}
By hypothesis, \(F_1'(X)\) is invertible in \(k(X)\); these two
equalities therefore imply
\[
  F_1(X)^{p^n}=1.
\]
Thus \(F_1\) is constant, contrary to the choice of \(F_1\).
Consequently \(F\) is constant and equal to \(1=F(0)\).
\qed
